% !PRINTED 23
\weberpage{23}
%<<CONT
und da nun $R$ ein grösster Normaltheiler von $Q$ und von $Q'$ ist,
so können wir daraus zwei neue Compositionsreihen von $P$ bilden,
nämlich
\weq{26}{\begin{aligned}
&P,\ Q,\ R,\ R_1,\ R_2,\ R_3 \ldots \\
&\hphantom{P,\ }\nu,\ \nu',\ \mu_1,\ \mu_2,\ \mu_3 \ldots
\end{aligned}}
\weq{27}{\begin{aligned}
&P,\ Q',\ R,\ R_1,\ R_2,\ R_3 \ldots \\
&\hphantom{P,\ }\nu',\ \nu,\ \mu_1,\ \mu_2,\ \mu_3 \ldots
\end{aligned}}
Die beiden Compositionsreihen (26) und (27) von $P$ haben also
dieselbe Indexreihe.

Nehmen wir jetzt an, der zu beweisende Satz sei bereits als richtig
erwiesen für Gruppen, deren Grad $\leqq \frac{1}{2}n$ ist (wenn $n$
den Grad von $P$ bedeutet), so folgt, dass alle Indexreihen von
$Q$, dessen Grad ja ein echter Theiler von $n$ und also höchstens
$= \frac{1}{2}n$ ist, mit einander übereinstimmen, dass also die
Reihen
\[
\nu',\ \mu_1,\ \mu_2,\ \mu_3 \ldots
\]
\[
\nu_1,\ \nu_2,\ \nu_3,\ \nu_4 \ldots,
\]
von der Anordnung abgesehen, übereinstimmen. Ebenso stimmen die
Indexreihen von $Q'$
\[
\nu,\ \mu_1,\ \mu_2,\ \mu_3 \ldots
\]
\[
\nu_1',\ \nu_2',\ \nu_3',\ \nu_4' \ldots
\]
überein. Also stimmen auch die beiden Indexreihen von $P$
\[
\nu,\ \nu_1,\ \nu_2,\ \nu_3,\ \nu_4 \ldots
\]
\[
\nu',\ \nu_1',\ \nu_2',\ \nu_3',\ \nu_4' \ldots
\]
mit einander überein, da die erste mit $\nu, \nu', \mu_1, \mu_2, \mu_3
\ldots$, die zweite mit $\nu', \nu, \mu_1, \mu_2, \mu_3 \ldots$
übereinstimmt.

Für Gruppen vom Grade 2, die nur den einen Normaltheiler 1 haben,
ist aber der Satz evident, und also ist er durch vollständige
Induction allgemein bewiesen.

Wenn in zwei Compositionsreihen von $P$, die wir mit ihren
Indexreihen jetzt so darstellen wollen
\weq{28}{\begin{aligned}
&P,\ P_1,\ P_2,\ P_3 \ldots P_{\mu-1},\ 1 \\
&\hphantom{P,\ }j_1,\ j_2,\ j_3 \ldots j_{\mu-1},\ j_\mu
\end{aligned}}
\weq{29}{\begin{aligned}
&P,\ P_1',\ P_2',\ P_3' \ldots P_{\mu-1}',\ 1 \\
&\hphantom{P,\ }j_1',\ j_2',\ j_3' \ldots j_{\mu-1}',\ j_\mu'
\end{aligned}}
ein gemeinschaftliches Glied $P_r = P_s'$ vorkommt, so gilt der
Satz von der Constanz der Indexreihen auch für die Gruppe $P_r =
P_s'$, und daraus folgt zunächst, dass $s = r$ sein muss, und
%CONT>>
